% !TeX root = ../../Book2.tex
% 纲要标题：### 14.2 Sard 定理的证明
% 纲要位置：几何学计划纲要.md，第 1549 行。
% 正文状态：初稿；尚未正式审核。

% 纲要内容开始
% 
% * 临界集合的秩分层
% * 秩为零部分的覆盖估计
% * 维数归纳与 Fubini 步骤
% * 正则性阶数及反例范围
% 
% 纲要内容结束

\section{Sard 定理的证明}
\label{g2:sec:1402}

方程 \(f(x)=y\) 的解集能否成为光滑子流形，取决于解上的微分是否满射。单独检查一个指定的 \(y\) 可能很难；Sard 定理给出整体信息：对光滑映射，坏的目标值只占一个零测集。证明必须控制临界点的\emph{像}，因为源中的临界点本身可能很多，甚至充满整个源流形。

\subsection{临界集合的秩分层}
\label{g2:subsec:1402:01}

先明确这里的“坏”。设 \(f:M^m\to N^n\) 是 \(C^1\) 映射，\(N\) 无边界。若 \(\mathrm{d}f_x:\mathrm{T}_xM\to\mathrm{T}_{f(x)}N\) 不满射，称 \(x\) 为\term{临界点}[critical point]；否则称为\term{正则点}[regular point]。临界点的像称为\term{临界值}[critical value]。若 \(f^{-1}(y)\) 中每一点都正则，称 \(y\) 为\term{正则值}[regular value]。特别地，不在像中的点也是正则值；一个值只要有一个临界原像，就已经是临界值。

\ProseParagraphBreak

在 \(n\) 维流形 \(N\) 中，集合 \(E\) 称为\term{零测集}[set of measure zero]，若它在每张坐标图中的像都有 \(n\) 维 Lebesgue 测度零。Euclid 空间中的含义是：给定任意 \(\varepsilon>0\)，能用可数个长方体覆盖它，并使体积总和小于 \(\varepsilon\)。这个条件与坐标无关。在较小的紧坐标块上，光滑坐标变换的导数有界，因而是 Lipschitz 映射；一个半径为 \(r\) 的小球的像落在半径至多为 \(Lr\) 的球中，覆盖总体积至多乘上固定常数。再用可数个这样的坐标块覆盖，便知零测性保持。这里用的是同维坐标变化，无须先在 \(N\) 上选择体积形式。

\begin{example}\label{g2:exm:sard-critical-values}
	映射 \(f:\mathbb R^2\to\mathbb R^2\)，\(f(x,y)=(x,y^2)\)，有
	\[
		\mathrm{D}f_{(x,y)}=\begin{pmatrix}1&0\\0&2y\end{pmatrix}.
	\]
	临界点集是 \(y=0\)，临界值集是目标中的横轴。第二坐标为正的值有两个正则原像，第二坐标为负的值没有原像；这两类值都正则。对于常值映射 \(M\to\mathbb R\)，整个 \(M\) 都由临界点组成，临界值却只有一个。
\end{example}

在坐标域 \(U\subset\mathbb R^m\) 内写 \(C=\{x:\operatorname{rank}\mathrm{D}f_x<n\}\)，并按秩分为
\[
	C^{(q)}=\{x:\operatorname{rank}\mathrm{D}f_x=q\},\qquad 0\le q<n.
\]
这些集合未必都是子流形。分层在这里首先是一种分类：当 \(q\ge1\) 时，某个非零的偏导数让我们把一个输出坐标直接改作输入坐标，进而把问题降低一维。

\begin{lemma}\label{g2:lem:sard-rank-reduction}
	设 \(f:U\to\mathbb R^n\) 光滑，\(U\subset\mathbb R^m\) 开，且 \(\operatorname{rank}\mathrm{D}f_a\ge1\)。重排坐标后，在 \(a\) 的某个邻域上存在光滑源坐标 \((t,z)\in J\times V\subset\mathbb R\times\mathbb R^{m-1}\)，使
	\[
		f(t,z)=(t,G(t,z)),\qquad
		\operatorname{rank}\mathrm{D}f_{(t,z)}=1+\operatorname{rank}\mathrm{D}_zG(t,z).
	\]
\end{lemma}
\begin{proof}
	选择非零偏导数并重排坐标，使 \(\partial f^1/\partial x^1(a)\ne0\)。令
	\[
		\Phi(x)=\bigl(f^1(x),x^2,\ldots,x^m\bigr).
	\]
	其导数行列式正是这个非零偏导数。第一卷第11章第3节的逆函数定理保证 \(\Phi\) 在充分小的邻域上可光滑反解；缩小像邻域，可取乘积盒 \(J\times V\)。在新坐标中，第一分量为 \(t\)，其余分量记作 \(G\)，于是
	\[
		\mathrm{D}f=\begin{pmatrix}1&0\\\partial_tG&\mathrm{D}_zG\end{pmatrix}.
	\]
	利用第一行消去第一列下方的元素，矩阵的秩便等于 \(1+\operatorname{rank}\mathrm{D}_zG\)。
\end{proof}

这里没有要求邻域内的秩恒定；非零偏导数只负责选出一个可用的坐标。真正需要继续处理的，正是其余方向上仍会变化的秩。

\subsection{秩为零部分的覆盖估计}
\label{g2:subsec:1402:02}

若 \(\mathrm{D}f_a=0\)，上面的坐标办法暂时失效。此时 Taylor 展开提供另一种控制：若若干阶导数都为零，一个边长为 \(\delta\) 的源立方体的像就比 \(\delta\) 小得多。源中小立方体的数量约为 \(\delta^{-m}\)，而每个像的 \(n\) 维体积约为 \(\delta^{n(k+1)}\)；当后一个指数更大时，总体积趋于零。

\begin{lemma}\label{g2:lem:sard-flat-cover}
	设 \(U\subset\mathbb R^m\) 开，\(m,n\ge1\)，\(f:U\to\mathbb R^n\) 为 \(C^{k+1}\) 映射，\(k\ge1\)。令
	\[
		A_k=\{x\in U:\partial^\alpha f^i(x)=0\text{，对所有 }1\le|\alpha|\le k,\ 1\le i\le n\}.
	\]
	若 \(n(k+1)>m\)，则 \(f(A_k)\) 为零测集。
\end{lemma}
\begin{proof}
	用可数个包含于 \(U\) 的闭立方体覆盖 \(U\)，只需在其中一个闭立方体 \(Q\) 上证明结论。\(Q\) 上各个 \(k+1\) 阶偏导数有共同的上界。将边长为 \(R\) 的 \(Q\) 均分成 \(s^m\) 个边长 \(R/s\) 的小立方体。对每个与 \(A_k\) 相交的小立方体 \(Q_j\)，选取 \(a_j\in A_k\cap Q_j\)。Taylor 公式的低阶项全部为零，故存在与 \(s,j\) 无关的常数 \(B\)，使
	\[
		|f(x)-f(a_j)|\le B|x-a_j|^{k+1}
		\le B\bigl(\sqrt m R/s\bigr)^{k+1}\qquad(x\in Q_j).
	\]
	因此 \(f(A_k\cap Q)\) 被至多 \(s^m\) 个半径为 \(B(\sqrt m R/s)^{k+1}\) 的球覆盖。若 \(\omega_n\) 表示单位 \(n\) 维球的体积，覆盖总体积至多为
	\[
		s^m\omega_n B^n(\sqrt m R/s)^{n(k+1)}
		=C_Qs^{m-n(k+1)}\longrightarrow0.
	\]
	零测性由此成立；可数并仍为零测集，故得到整个 \(A_k\) 的结论。
\end{proof}

这里的多重指标 \(\alpha=(\alpha_1,\ldots,\alpha_m)\) 记录各坐标求导次数，\(|\alpha|=\alpha_1+\cdots+\alpha_m\)。上式只需普通的有限阶 Taylor 余项估计，不要求函数等于自己的 Taylor 级数。对于光滑函数，在一点的全部导数都为零也不能推出其在邻域内恒定；例如 \(e^{-1/x^2}\) 在 \(x=0\) 补值为零后的函数就具有这种现象。覆盖估计恰好能处理这些极其平坦的点。

\subsection{维数归纳与 Fubini 步骤}
\label{g2:subsec:1402:03}

剩下的工作是把前面两种控制接起来。正秩部分沿一个输出坐标切片；零秩但尚未高度平坦的部分，则落在某个偏导数的正则零集上。两种情形都把源的维数减少一。

\ProseParagraphBreak

切片时使用如下 Fubini 结论：若 \(E\subset\mathbb R\times\mathbb R^{n-1}\) 紧，\(n\ge2\)，且每个截面 \(E_t=\{w:(t,w)\in E\}\) 都为 \((n-1)\) 维零测集，则 \(E\) 为 \(n\) 维零测集。因为紧集是 Borel 可测集，对其示性函数使用非负函数的 Fubini--Tonelli 定理，得到
\[
	\operatorname{vol}_n(E)=\int_{\mathbb R}\operatorname{vol}_{n-1}(E_t)\,\mathrm{d}t=0.
\]
这一紧集切片结论也有只用有限覆盖的证明，见\cite[第126页，引理6.2]{Lee2012Smooth}。下面每次都先限制到紧块，使所处理的像确实可测。

\begin{theorem}[Sard 定理]\label{g2:thm:sard}
	设 \(M^m\) 是有限维、Hausdorff、第二可数的光滑流形，允许有边界；\(N^n\) 是无边界光滑流形，具有相同的分离性与可数性条件。若 \(f:M\to N\) 光滑，则它的临界值集在 \(N\) 中为零测集。因此正则值在 \(N\) 中稠密。
\end{theorem}
\begin{proof}
	先证明无边界情形，并对源维数 \(m\) 归纳。当 \(n=0\) 时，所有微分都满射，临界值集为空。当 \(m=0<n\) 时，第二可数的零维流形是可数离散集，像也可数，故为零测集。以下设 \(m,n\ge1\)，并假设所有源维数小于 \(m\) 的光滑映射均已满足结论。

	取可数的源、目标坐标图，并将源图限制在各目标图的原像内。坐标变化保持微分的秩和目标零测性，故只需讨论 \(f:U\subset\mathbb R^m\to\mathbb R^n\)。除临界集 \(C\) 外，记 \(A_j\) 为前一引理中一至 \(j\) 阶偏导数全为零的集合。它们在 \(U\) 内闭，并满足
	\[
		C\supset A_1\supset A_2\supset\cdots.
	\]
	取整数 \(k\ge1\) 使 \(n(k+1)>m\)。我们依次处理
	\[
		C=(C\setminus A_1)\cup\bigcup_{j=1}^{k-1}(A_j\setminus A_{j+1})\cup A_k.
	\]

	对 \(C\setminus A_1\)，若 \(n=1\)，此集合为空：一维目标的微分不满射就只能为零。若 \(n\ge2\)，\cref{g2:lem:sard-rank-reduction}在每个这样的点附近给出 \(f(t,z)=(t,G(t,z))\)。令 \(G_t(z)=G(t,z)\)。秩公式说明 \((t,z)\) 是 \(f\) 的临界点，当且仅当 \(z\) 是 \(G_t:V\to\mathbb R^{n-1}\) 的临界点。对源坐标盒内任意紧块 \(K\)，集合 \(E=f(C\cap K)\) 紧，而它在目标超平面第一坐标等于 \(t\) 上的截面包含于 \(G_t\) 的临界值集。\(G_t\) 的源维数是 \(m-1\)，由归纳假设，每个这样的截面为零测集；Fubini 结论遂给出 \(E\) 零测。用可数紧块穷尽各坐标盒，再用可数个坐标盒覆盖 \(C\setminus A_1\)，得到 \(f(C\setminus A_1)\) 零测。

	对 \(A_j\setminus A_{j+1}\)，取一点 \(a\)。某个分量的某个 \(j+1\) 阶偏导数在 \(a\) 非零，所以能选一个 \(j\) 阶偏导数 \(h=\partial^\alpha f^i\)，使 \(\mathrm{D}h_a\ne0\)。同时 \(h\) 在 \(A_j\) 上为零。第一卷第11章第3节的隐函数定理说明，在 \(a\) 的充分小邻域 \(V_a\) 内，\(H_a=\{h=0\}\) 是光滑超曲面：选取 \(h\) 的一个非零偏导数，即可把相应坐标解成其余 \(m-1\) 个坐标的光滑函数。于是 \(A_j\cap V_a\subset H_a\)。在这些点上，\(\mathrm{d}f=0\)，故限制映射 \(f|_{H_a}\) 的微分也为零，特别地不满射。归纳假设应用于源维数为 \(m-1\) 的 \(f|_{H_a}\)，说明 \(f(A_j\cap V_a)\) 零测。用可数个这样的邻域覆盖 \(A_j\setminus A_{j+1}\)，便得到这一整部分的像零测。

	最后，\cref{g2:lem:sard-flat-cover}直接给出 \(f(A_k)\) 零测。三部分合起来完成归纳。

	若 \(M\) 有边界，在每张边界坐标图中，光滑映射的坐标表达按定义能在每一点附近延拓到半空间之外的开邻域。原映射在边界点的微分就是这个局部延拓的微分；所以原映射的临界点也是延拓映射的临界点。对可数个局部延拓使用刚证的无边界结论，仍得到临界值集零测。最后，当 \(n\ge1\) 时，任意非空目标开集在某张坐标图中包含一个小球，具有正测度，不能全部由临界值组成；\(n=0\) 时所有值都正则。因此正则值稠密。
\end{proof}

这条证明的分工可由下表记忆。它按照点上的导数选择工具，而没有把整片临界集误当成一个光滑子流形。

\begin{center}
	\begin{tabularx}{\textwidth}{@{}lXX@{}}
		\toprule
		临界点所在部分 & 可用的信息 & 控制临界值的方法 \\
		\midrule
		\(C\setminus A_1\) & 某个一阶偏导数非零 & 一个输出坐标作参数，降维后用 Fubini \\
		\(A_j\setminus A_{j+1}\) & 某个高阶偏导数有正则零集 & 限制到包含该部分的超曲面，降低源维数 \\
		\(A_k\) & 低阶 Taylor 项全部消失 & 小立方体数量与像球体积的覆盖估计 \\
		\bottomrule
	\end{tabularx}
\end{center}

\begin{corollary}\label{g2:cor:sard-lower-dimension}
	设 \(f:M^m\to N^n\) 是上述流形间的光滑映射，且 \(m<n\)。则 \(f(M)\) 是 \(N\) 中的零测集；每个正则值的原像均为空。
\end{corollary}
\begin{proof}
	从 \(m\) 维空间到 \(n\) 维空间的线性映射不可能满射，因此每一点都是临界点。像恰为临界值集，应用\cref{g2:thm:sard}即可。
\end{proof}

例如光滑曲线在平面中的像具有二维测度零，即使曲线在非紧参数区间上的像可能稠密，仍不能因此获得正面积。“稠密”说的是闭包，“零测”说的是体积，二者衡量不同的性质。

\subsection{正则性阶数及反例范围}
\label{g2:subsec:1402:04}

光滑性在上述归纳中使任意阶偏导数都继续光滑。对于有限阶可微映射，\(h=\partial^\alpha f^i\) 的可微阶数会随 \(|\alpha|\) 增大而降低，隐函数定理所得超曲面的正则性也随之降低。因此，那个归纳过程不能直接给出最优的有限阶条件。

\ProseParagraphBreak

有限阶 Sard 定理的准确形式是：对 \(C^r\) 映射 \(f:M^m\to N^n\)，若整数 \(r\ge1\) 满足
\[
	r>\max\{m-n,0\},
\]
则临界值集为零测集。这里 \(M,N\) 可取光滑流形，映射只要求 \(C^r\)，源的边界按局部延拓处理。有限阶版本的证明参见\cite[第885--889页，定理4.1、7.2]{Sard1942}；其中高阶平坦集合的细分估计负责避免上述正则性损失。本章后续参数定理使用已完整证明的光滑版本。

\ProseParagraphBreak

这个阶数条件有可见的几何意义。对平面上的函数，\(C^2\) 足够；单有 \(C^1\) 时，临界值甚至可以覆盖整个区间。下面给出一个具体构造，使“更多可微性”所排除的现象看得见。

\begin{example}\label{g2:exm:sard-c1-counterexample}
	存在 \(C^1\) 函数 \(F:\mathbb R^2\to\mathbb R\)，其临界值集包含 \([0,2]\)。

	从 \([0,1]\) 开始，每一步保留每个区间两端各占其长度 \(\lambda=2/5\) 的闭子区间，得到 Cantor 集 \(E\)。第 \(k\) 步删去 \(2^{k-1}\) 个开间隙，每个长度为 \((1-2\lambda)\lambda^{k-1}\)。固定一个在 \((0,1)\) 上为正、端点处平坦、积分为 \(1\) 的光滑函数 \(\rho\)，并在每个第 \(k\) 步间隙 \(I=(a,b)\) 上放置
	\[
		h(x)=\frac{3^{-k}}{b-a}\rho\left(\frac{x-a}{b-a}\right).
	\]
	在 \(E\) 及 \([0,1]\) 外令 \(h=0\)。间隙上这项的积分为 \(3^{-k}\)，上确界至多为
	\[
		\frac{\|\rho\|_\infty}{3(1-2\lambda)}
		\left(\frac1{3\lambda}\right)^{k-1}\longrightarrow0.
	\]
	这保证 \(h\) 连续：在某个固定间隙的端点附近，连续性来自 \(\rho\) 在端点趋于零；当间隙的级数趋于无穷时，则由上述一致上界趋于零。有限多个低级间隙与余下间隙分别处理，即得到每个 \(E\) 中点处的连续性。

	令 \(g(x)=\int_0^xh(t)\,\mathrm{d}t\)，则 \(g\in C^1(\mathbb R)\)，且 \(g'=0\) 于 \(E\)。全部间隙的积分之和为 \(\sum_{k\ge1}2^{k-1}3^{-k}=1\)。一个第 \(k\) 步保留下来的基本区间，其内各级间隙的积分之和为
	\[
		\sum_{j=k+1}^{\infty}2^{j-k-1}3^{-j}=3^{-k}.
	\]
	所以 \(g\) 把保留区间依次映到标准三分 Cantor 构造中的保留区间，且把每个间隙映到相应的中间三分之一。于是 \(g(E)=C\)，其中 \(C\) 是标准三分 Cantor 集。

	现在取 \(F(x,y)=g(x)+g(y)\)。在 \(E\times E\) 上，\(\mathrm{d}F=(g'(x),g'(y))=0\)，而 \(F(E\times E)=C+C=[0,2]\)。最后这个等式可这样验证：\([-1,1]\) 中每个数都能写成 \(\sum_{j\ge1}2\varepsilon_j3^{-j}\)，其中 \(\varepsilon_j\in\{-1,0,1\}\)，因为第一位所给的三个闭区间覆盖 \([-1,1]\)，选定第一位后再缩放余项即可递归选择。把正、负位分别放到两个只含 \(0,2\) 的三进制展开里，得到 \(C-C=[-1,1]\)。再由 \(1-C=C\)，得到 \(C+C=1+(C-C)=[0,2]\)。
\end{example}

因此，零测结论涉及映射的正则性，而不只是处处存在一个连续的微分。在一般维数中，整数阶阈值的最优性及相应构造见\cite[第890页，第8节]{Sard1942}；对当前 \(m=2,n=1\) 的情形，上例已经完整展示 \(C^1\) 与 \(C^2\) 的差别。另一个边界是源和目标维数：当 \(m<n\) 时，\(C^1\) 已足以让整个像零测：在一个紧坐标立方体上，导数有界使映射 Lipschitz；分割所得的 \(s^m\) 个小立方体各自的像都被半径至多 \(B/s\) 的球包含，故总体积至多为 \(Cs^{m-n}\to0\)。再取可数坐标覆盖即可。
